
本小节做的实验是为了把通过实验1验证的代码部署到实际的FPGA开发板上，即实验目的为： 熟悉使用 Vivado 进行硬件开发的步骤，了解综合、布线、在线调试的全过程。





\subsection{实验原理}
\subsubsection{ILA测试}
ILA可以用来观测FPGA内部信号，其核心在于利用一个定时器来截获特定的信号，
即采样，当采样频率不够（未高于采样频率的2倍），则会产生失真，或者说我们采样到信号波形不能充分表征原始的信号。相反，采样频率升高的时候，采样观察边沿的时间分辨率就可以提高。
但如果只是提高采样频率，而不提高记录的时长，最后也可能导致想看的波形被截断而不能得到充分的观察，
我们需要去同时考虑采样点数的失真效应和记录时长带来的截断效应。
采样点数$ N_{sample} $，采样间隔$ T_{inter} $和记录时长$ T_{all} $有着如下的关系：
\[
\frac{N_{sample}}{T_{inter}} = T_{all}
\]

而在本次实验中，ILA的使用步骤为
\begin{enumerate}
    \item 在 IP Catalog 中搜索并添加 ILA IP，设置采样时钟与探针数量、位宽（Probes）；
    \item 在顶层将待观测信号（\texttt{clk\_100M}、\texttt{clk\_20M}、texttt{clk\_10M}、texttt{clk\_1k}）接入 ILA 探针，采样时钟选用 200MHz；
    \item 重新综合、实现并生成比特流，按上板流程下载；
    \item 在 Hardware Manager 中设置触发条件（如信号上升沿），点击运行按钮捕获波形。
\end{enumerate}


\subsection{实验过程}
实验步骤如下：
\begin{enumerate}
    \item 完成综合、实现后，点击 Generate Bitstream 生成比特流文件（.bit）；
    \item 用 USB 下载器连接实验箱 JTAG 接口，接通电源（注意单排针方向，勿带电插拔）；
    \item 打开 Hardware Manager → Open Target → Auto Connect 连接 FPGA；
    \item 右键器件 → Program Device，选择 .bit 文件完成下载；
    \item 下载成功后观察板上现象，配合 ILA 波形验证各路时钟是否正确。
\end{enumerate}

预期可以观察到：用ILA去观测产生的100MHz，20MHz，10MHz，1kHz的信号，其频率正确并且占空比为预期的50\%


一些实验过程图展示：



\subsection{实验结果}
